% !TEX root = ../main.tex
\chapter{ケーススタディ：User スキーマの lossless migration}
\label{chap:ch28_user_schema_lossless_migration}
\BookIndex{lossless migration}{lossless migration}
\BookIndex{すきーまいこう}{スキーマ移行}
\BookIndex{どうけい}{同型}

\begin{chapterabstract}
本章では、\texttt{UserV1} から \texttt{UserV2} へのスキーマ変更を、情報を失わない移行として扱います。
フィールド名の変更、ネスト構造の導入、関連フィールドのグルーピングは、見た目には大きな変更に見えます。
しかし、旧形式から新形式へ移し、さらに旧形式へ戻したときに元の値へ戻るなら、その範囲では情報は失われていません。
逆向きにも、新形式から旧形式へ戻してから新形式へ戻したときに元の値へ戻るなら、二つのスキーマは同じ情報を別の形で表しているといえます。
この性質を Lean 4 では round-trip law として書き、機械的に検査します。
\end{chapterabstract}

\begin{learninggoals}
\item \texttt{UserV1} と \texttt{UserV2} の違いを、情報の違いと表現の違いに分けて説明できます。
\item \texttt{migrate : UserV1 -> UserV2} と \texttt{rollback : UserV2 -> UserV1} を Lean で定義できます。
\item \texttt{rollback (migrate u) = u} と \texttt{migrate (rollback v) = v} を round-trip 性質として読めます。
\item lossless migration と、単なる片方向変換やデフォルト値付き移行を区別できます。
\item Lean の証明を CI や設計レビューに入れるときの境界線を説明できます。
\end{learninggoals}

\section{ユーザースキーマを情報損失なく移行する}

ユーザー情報のスキーマは、長く運用されるサービスほど変わりやすくなります。
最初は単純なフラットなレコードだったものが、後からプロフィール、連絡先、監査情報、権限情報などに分かれます。
API のレスポンス形式も、初期は小さく始まり、後から用途ごとのまとまりを持ちます。

このような変更では、実装者はよく「同じ情報を移しているだけ」と述べます。
しかし、その言葉だけでは安全性を確認できません。
どのフィールドがどこへ移ったのでしょうか。
名前を変えただけなのでしょうか。
複数の値をまとめただけなのでしょうか。
新しい値を補っていないでしょうか。
どこかで正規化や切り捨てをしていないでしょうか。
これらを明示しないと、移行レビューは雰囲気に頼ることになります。

図\ref{fig:ch28-lossless-migration-overview}は、旧形式と新形式の各フィールドの対応、および両方向の変換を示します。

\FigurePlaceholder{fig-ch28-lossless-migration-overview.png}{0.90}{User スキーマの相互変換}{fig:ch28-lossless-migration-overview}

\section{ソフトウェア上の問題}

本章の例では、旧スキーマ \texttt{UserV1} を次のようなフラットな構造として考えます。

\begin{itemize}
\item \texttt{id}: ユーザーの識別子
\item \texttt{name}: 表示名
\item \texttt{email}: メールアドレス
\item \texttt{emailVerified}: メール確認済みかどうか
\item \texttt{createdAt}: 作成時刻を表す簡略化した数値
\end{itemize}

新スキーマ \texttt{UserV2} では、情報のまとまりを明示します。
識別子は \texttt{userId} へ名前を変えます。
表示名は \texttt{Profile} に入れます。
メールアドレスと確認状態は \texttt{Contact} に入れます。
作成時刻は \texttt{Audit} に入れます。

\begin{table}[htbp]
\centering
\caption{フィールド対応}
\label{tab:ch28-field-mapping}
\begin{tabular}{lll}
\toprule
旧スキーマ & 新スキーマ & 変更の種類 \\
\midrule
\texttt{id} & \texttt{userId} & フィールド名変更 \\
\texttt{name} & \texttt{profile.displayName} & ネスト構造への移動 \\
\texttt{email} & \texttt{contact.emailAddress} & ネスト構造への移動と名前変更 \\
\texttt{emailVerified} & \texttt{contact.verified} & 関連値のグルーピング \\
\texttt{createdAt} & \texttt{audit.createdAt} & 監査情報へのグルーピング \\
\bottomrule
\end{tabular}
\end{table}

\Cref{tab:ch28-field-mapping} は、各フィールドがどこへ移ったかを示しています。
表の右端に注目すると、今回の変更は、分割、結合、正規化、削除ではありません。
すべての値が一対一に対応しています。
これが、lossless migration として扱うための出発点です。

\section{圏論での読み方：同型としてのスキーマ変更}

圏論では、情報を失わない相互変換を同型として見ます。
本章では、対象をデータ型、射を変換関数として読みます。
すると、\texttt{UserV1} と \texttt{UserV2} が同型であるとは、次の四つがあることを意味します。

\textbf{スキーマ同型}とは、二つのスキーマ \(A\) と \(B\) について、次のデータがあることです。
\begin{itemize}
\item \(A\) から \(B\) への変換 \(\mathit{to}\)
\item \(B\) から \(A\) への変換 \(\mathit{from}\)
\item 任意の \(a : A\) について \(\mathit{from}(\mathit{to}(a)) = a\)
\item 任意の \(b : B\) について \(\mathit{to}(\mathit{from}(b)) = b\)
\end{itemize}
本章の名前では、\(\mathit{to}\) が \texttt{migrate}、\(\mathit{from}\) が \texttt{rollback} に対応します。

図\ref{fig:ch28-round-trip-laws}の左側は旧形式上、右側は新形式上で合成が恒等変換になる法則を表します。

\FigurePlaceholder{fig-ch28-round-trip-laws.png}{0.90}{round-trip 法則}{fig:ch28-round-trip-laws}

ここで大切なのは、変換関数が二つ存在するだけでは不十分だという点です。
\texttt{migrate} と \texttt{rollback} を実装できても、それらが互いの逆になっていなければ同型ではありません。
実務では「戻す関数もあります」と述べるだけで安全に見えることがあります。
しかし、戻した結果が常に元の値と一致することを確認して初めて、lossless といえます。

\section{Lean 4 でスキーマを定義する}

まず、Lean 上に小さなモデルを作ります。
現実のユーザー情報には、UUID、メール形式、タイムゾーン、監査ログ、外部認証 ID などが関わります。
本章では、round-trip の核心を見やすくするため、識別子や時刻を \texttt{Nat}、文字列を \texttt{String}、確認状態を \texttt{Bool} で表します。

\begin{listing}[htbp]
\begin{minted}{lean4}
namespace Chapter28

structure UserV1 where
  id : Nat
  name : String
  email : String
  emailVerified : Bool
  createdAt : Nat
  deriving Repr, DecidableEq

structure Profile where
  displayName : String
  deriving Repr, DecidableEq

structure Contact where
  emailAddress : String
  verified : Bool
  deriving Repr, DecidableEq

structure Audit where
  createdAt : Nat
  deriving Repr, DecidableEq

structure UserV2 where
  userId : Nat
  profile : Profile
  contact : Contact
  audit : Audit
  deriving Repr, DecidableEq
\end{minted}
\caption{\texttt{UserV1} と \texttt{UserV2}}
\label{lst:ch28-schema-defs}
\end{listing}

\texttt{deriving Repr} は、\texttt{\#eval} で値を表示しやすくするために付けています。
\texttt{deriving DecidableEq} は、値の等しさを計算で確認したい場面にも備えるためです。
ただし、本章の中心は、テスト的な比較ではなく、すべての入力に対する定理です。

\section{\texttt{migrate} を定義する}

次に、旧スキーマから新スキーマへ移す関数を定義します。
ここでは、値を加工しません。
文字列を trim しません。
メールアドレスを小文字化しません。
時刻を変換しません。
単に、対応する場所へ値を置くだけです。

\begin{listing}[htbp]
\begin{minted}{lean4}
def migrate (u : UserV1) : UserV2 :=
  { userId := u.id
    profile := { displayName := u.name }
    contact :=
      { emailAddress := u.email
        verified := u.emailVerified }
    audit := { createdAt := u.createdAt } }

#eval migrate
  { id := 1
    name := "Ada"
    email := "ada@example.com"
    emailVerified := true
    createdAt := 1000 }
-- =>
-- { userId := 1,
--   profile := { displayName := "Ada" },
--   contact := { emailAddress := "ada@example.com", verified := true },
--   audit := { createdAt := 1000 } }
\end{minted}
\caption{\texttt{migrate}}
\label{lst:ch28-migrate}
\end{listing}

この関数の型は \texttt{UserV1 -> UserV2} です。
ソフトウェアの言葉では、旧 DTO から新 DTO への adapter と読めます。
圏論の言葉では、\texttt{UserV1} という対象から \texttt{UserV2} という対象への射として読めます。

この段階では、まだ lossless だとはいえません。
\texttt{migrate} は、片方向の変換にすぎません。
値をすべて移しているように見えても、逆向きに戻して本当に一致することをまだ証明していません。

\section{\texttt{rollback} を定義する}

逆向きの変換も定義します。
\texttt{rollback} は、新スキーマの各ネスト構造から値を取り出し、旧スキーマのフラットなフィールドへ戻します。

\begin{listing}[htbp]
\begin{minted}{lean4}
def rollback (v : UserV2) : UserV1 :=
  { id := v.userId
    name := v.profile.displayName
    email := v.contact.emailAddress
    emailVerified := v.contact.verified
    createdAt := v.audit.createdAt }
\end{minted}
\caption{\texttt{rollback}}
\label{lst:ch28-rollback}
\end{listing}

\texttt{rollback} が存在するだけなら、単に \texttt{UserV2 -> UserV1} の関数があるというだけです。
たとえば、\texttt{name} に常に空文字を入れるような \texttt{rollback} も、型だけを見れば定義できます。
したがって、次に必要なのは round-trip の証明です。

\section{証明する性質を選ぶ}

このケースで証明したい性質は、二つあります。

\begin{align*}
\texttt{rollback (migrate u)} &= \texttt{u} \\
\texttt{migrate (rollback v)} &= \texttt{v}
\end{align*}

一つ目は、旧データを新形式へ移し、旧形式へ戻すと、元の旧データに一致することを述べます。
これは、既存データの移行にとって特に重要です。
二つ目は、新データを旧形式へ戻し、再び新形式へ移すと、元の新データに一致することを述べます。
これは、新スキーマ全体が旧スキーマと同じ情報だけからできていることを表します。

\textbf{両方向 round-trip} は、片方だけの互換性より強い性質です。
\texttt{rollback (migrate u) = u} だけなら、旧データを安全に新形式へ移せることを示します。
しかし、\texttt{migrate (rollback v) = v} も成り立つとき、新スキーマ側にも余分な情報や戻せない自由度がないことを示せます。

\section{Lean で round-trip を証明する}

証明方針は単純です。
\texttt{UserV1} の値 \texttt{u} をフィールドに分解します。
すると、\texttt{migrate} は各フィールドを対応先へ置き、\texttt{rollback} はそれを元の場所へ戻します。
最終的に左右の式は同じ構造体リテラルになるので、\texttt{rfl} で閉じられます。

\begin{listing}[htbp]
\begin{minted}{lean4}
theorem rollback_migrate (u : UserV1) :
    rollback (migrate u) = u := by
  cases u
  rfl
\end{minted}
\caption{\texttt{rollback\_migrate}}
\label{lst:ch28-rollback-migrate}
\end{listing}

\texttt{cases u} は、\texttt{u} をそのフィールドに分解します。
分解後、Lean は \texttt{migrate} と \texttt{rollback} の定義を展開できます。
左右が同じ形になったので、\texttt{rfl}、つまり反射律で証明できます。

新から旧へ戻して再び新へ移す証明では、\texttt{UserV2} の中に \texttt{Profile}、\texttt{Contact}、\texttt{Audit} が入っています。
そのため、外側の \texttt{UserV2} だけでなく、内側の構造体も分解します。

\begin{listing}[htbp]
\begin{minted}{lean4}
theorem migrate_rollback (v : UserV2) :
    migrate (rollback v) = v := by
  cases v with
  | mk userId profile contact audit =>
    cases profile with
    | mk displayName =>
      cases contact with
      | mk emailAddress verified =>
        cases audit with
        | mk createdAt =>
          rfl
\end{minted}
\caption{\texttt{migrate\_rollback}}
\label{lst:ch28-migrate-rollback}
\end{listing}

この証明が少し長く見えるのは、新スキーマがネストしているためです。
ただし、行っていることは難しくありません。
各構造体をフィールドへ展開し、定義を計算した結果、同じ値に戻ることを Lean に確認させています。

\section{同型を小さな構造体としてまとめる}

二つの round-trip 定理がそろったら、それらを一つの値としてまとめられます。
これは、\ChapterRef{chap:ch12-isomorphism}で扱った「同型を構造体として持つ」考え方の実務版です。

\begin{listing}[htbp]
\begin{minted}{lean4}
structure SchemaIso (A B : Type) where
  to : A -> B
  backward : B -> A
  backward_to : forall a, backward (to a) = a
  to_backward : forall b, to (backward b) = b

def userSchemaIso : SchemaIso UserV1 UserV2 :=
  { to := migrate
    backward := rollback
    backward_to := rollback_migrate
    to_backward := migrate_rollback }
\end{minted}
\caption{\texttt{SchemaIso}}
\label{lst:ch28-schema-iso}
\end{listing}

\texttt{userSchemaIso} は、単なる設計メモではありません。
\texttt{migrate} と \texttt{rollback} に加えて、二つの法則そのものを持つ値です。
この形にしておくと、「この変換は同型として扱える」という主張を、関数と証明の組として受け渡せます。

\section{lossless でなくなる典型例}

本章の例は、意図的に lossless になるように設計しています。
しかし、実務では似た見た目でも lossless ではない変更が多くあります。
たとえば、次の変更には注意が必要です。

\begin{itemize}
\item \texttt{name} を \texttt{firstName} と \texttt{lastName} に分けます。
\item \texttt{email} を常に小文字化して保存します。
\item \texttt{createdAt} を日単位へ丸めます。
\item \texttt{UserV2} に \texttt{timezone} のような新しい必須フィールドを追加します。
\item 旧スキーマの値を削除し、新スキーマではデフォルト値で補います。
\end{itemize}

これらは戻せる場合もありますが、常に戻せるとは限りません。
たとえば、\texttt{timezone} を新しく追加し、旧データからは常に \texttt{"UTC"} を入れるとします。
このとき、旧データから新データへの移行はできます。
しかし、任意の \texttt{UserV2} が \texttt{timezone = "UTC"} とは限りません。
そのため、\texttt{migrate (rollback v) = v} は一般には成り立ちません。

「旧データを新データへ移せる」ことと、「新旧スキーマが同じ情報を表す」ことは違います。
前者は片方向互換であり、後者は両方向 round-trip を必要とします。
本章の Lean 証明は、この二つを混同しないための安全装置です。

\section{CI に入れる場合の境界線}

Lean の証明を CI に入れると、スキーマ移行の設計変更に対するセンサーとして使えます。
たとえば、\texttt{migrate} の対応先フィールドを間違えたり、\texttt{rollback} で別の値を返したりすると、round-trip 定理が壊れます。
この失敗は、テストの失敗と同じようにレビューのきっかけになります。

図\ref{fig:ch28-ci-boundary}では、CI が検査する型、変換、round-trip と、JSON、ORM、DB 制約、運用手順など別途検査する層を分けます。

\FigurePlaceholder{fig-ch28-ci-boundary.png}{0.90}{CI と実システムの境界}{fig:ch28-ci-boundary}

ただし、CI に入れた Lean 証明が保証するのは、Lean でモデル化した範囲です。
現実のシステムには、JSON encoder、ORM、DB 制約、NULL、文字コード、時刻フォーマット、外部 API、並行実行、ロールバック手順などがあります。
本章の証明だけで、それらすべてが正しいとはいえません。

Lean はフィールド対応と round-trip を確認します。
ユニットテストは encoder、decoder、DB mapper との対応を確認します。
結合テストやリハーサルは、実 DB と運用手順を確認します。

CI に置くファイル名は、たとえば \texttt{lean/ch28\_user\_schema\_lossless\_migration.lean} とします。
定理名は、仕様項目に対応させるとよいでしょう。
\texttt{rollback\_migrate} は、「旧データを移して戻すと元に戻る」という仕様に対応します。
\texttt{migrate\_rollback} は、「新データを戻して再移行すると元に戻る」という仕様に対応します。
レビューでは、この二つの定理が何を保証し、何を保証しないかを明記します。

\section{本章のまとめ}
\enlargethispage{2\baselineskip}

lossless migration は、旧形式と新形式が同じ情報を別の構造で持つ変更です。

\texttt{migrate} と \texttt{rollback} があるだけでは同型ではありません。

\texttt{rollback (migrate u) = u} と \texttt{migrate (rollback v) = v} の二つを証明して初めて、両方向の情報保存を主張できます。

Lean では、構造体、関数、\texttt{cases}、\texttt{rfl} だけで小さな round-trip 証明を書けます。

CI に入れる場合は、Lean が保証するモデル上の性質と、本番システムで別途確認すべき性質を分けて書く必要があります。
